```latex
\documentclass[11pt,a4paper]{article}

\usepackage{amsmath,amssymb,amsthm,mathtools}
\usepackage[margin=1in]{geometry}
\usepackage{microtype}

\newtheorem{theorem}{Theorem}[section]
\newtheorem{lemma}[theorem]{Lemma}
\newtheorem{proposition}[theorem]{Proposition}
\newtheorem{corollary}[theorem]{Corollary}
\theoremstyle{remark}
\newtheorem{remark}[theorem]{Remark}

\numberwithin{equation}{section}

\newcommand{\R}{\mathbb{R}}
\newcommand{\E}{\mathbb{E}}
\newcommand{\Prob}{\mathbb{P}}
\newcommand{\W}{\mathbb{W}}
\newcommand{\bone}{\mathbf 1}
\newcommand{\dist}{\operatorname{dist}}
\newcommand{\supp}{\operatorname{supp}}
\newcommand{\sgn}{\operatorname{sgn}}
\newcommand{\dd}{\mathrm d}
\newcommand{\Div}{\operatorname{div}}
\newcommand{\KL}{\operatorname{KL}}
\newcommand{\Fcal}{\mathcal F}
\newcommand{\Ccal}{\mathcal C}
\newcommand{\Ncal}{\mathcal N}
\newcommand{\Scal}{\mathcal S}
\newcommand{\Tcal}{\mathcal T}
\newcommand{\norm}[1]{\left\|#1\right\|}
\newcommand{\inner}[2]{\left\langle #1,#2\right\rangle}
\newcommand{\abs}[1]{\left|#1\right|}

\title{Minimax estimation of the conductivity in Darcy's problem}
\date{}

\begin{document}

\maketitle

\begin{abstract}
We determine the sharp logarithmic minimax rate for estimating a
H\"older--Zygmund conductivity in the divergence-form Darcy equation on the
unit ball from indirect Gaussian white-noise observations of the potential.
The characteristic set of the linearised inverse problem is the
codimension-one set of covectors orthogonal to the radial vector field
$x/|x|$. Consequently, the effective number of severely ill-posed modes at
frequency $2^J$ is $2^{(d-1)J}$, rather than $2^{dJ}$. The sharp squared
$L^2$-risk is therefore
\[
n^{-\frac{2\alpha}{2\alpha+d+3}}
\]
up to logarithmic factors. An isotropic second-order regularisation argument
would give the non-sharp exponent $2\alpha/(2\alpha+d+4)$.
\end{abstract}

\section{Statement of the result}

Let $D=\{x\in\R^d:|x|<1\}$, $d\geq 2$, and let
$C^\alpha(D)=B^\alpha_{\infty,\infty}(D)$ denote the H\"older--Zygmund
space. For $\alpha>1$ and $0<r<1$, set
\[
\Fcal(\alpha,r)
=
\left\{
f\in C^\alpha(D):\|f-\bone\|_{C^\alpha(D)}\leq r
\right\}.
\]
Since the H\"older norm dominates the supremum norm,
\[
1-r\leq f(x)\leq 1+r,
\]
so every $f\in\Fcal(\alpha,r)$ is uniformly elliptic for $r<1$.

For $f\in\Fcal(\alpha,r)$, let $u_f$ be the unique weak solution of
\begin{equation}\label{eq:Darcy}
\Div(f\nabla u_f)=1\quad\text{in }D,
\qquad
u_f=0\quad\text{on }\partial D.
\end{equation}
Let $\W$ be Gaussian white noise on $L^2(D)$, and observe
\begin{equation}\label{eq:model}
Y=u_f+\frac{1}{\sqrt n}\W.
\end{equation}
Equivalently, for every $\varphi\in L^2(D)$,
\[
\langle Y,\varphi\rangle
=
\langle u_f,\varphi\rangle+\frac1{\sqrt n}W(\varphi),
\qquad
W(\varphi)\sim N(0,\|\varphi\|_{L^2}^2).
\]

The main result is as follows.

\begin{theorem}\label{thm:main}
Let
\[
\gamma=2\alpha+d+3,
\qquad
\kappa=\frac{2\alpha}{2\alpha+d+3}.
\]
There exist $r_0\in(0,1)$ and constants $c,C>0$, depending only on
$d,\alpha$, and $r$, such that for every fixed $0<r\leq r_0$ and all
$n\geq 2$,
\begin{equation}\label{eq:rates}
c n^{-\kappa}
\leq
\inf_{\widehat f_n}
\sup_{f_0\in\Fcal(\alpha,r)}
\E_{f_0}\!\left[
\|\widehat f_n-f_0\|_{L^2(D)}^2
\right]
\leq
C n^{-\kappa}(\log n)^\kappa .
\end{equation}
Thus, up to logarithmic factors,
\[
\boxed{\;
\inf_{\widehat f_n}\sup_{f_0\in\Fcal(\alpha,r)}
\E_{f_0}\|\widehat f_n-f_0\|_{L^2(D)}^2
\asymp
n^{-\frac{2\alpha}{2\alpha+d+3}} .
\;}
\]
\end{theorem}

The additive constant $3$ in $\gamma$ has the following origin. The inverse
of the divergence-form elliptic operator gains two derivatives, but the
derivative of the conductivity-to-solution map is severely ill-posed only in
directions tangential to the radial field $\nabla u_{\bone}(x)=x/d$. At
spatial frequency $2^J$, those tangential modes have multiplicity
$2^{(d-1)J}$. Hence the variance term has order
$2^{(d+3)J}/n$ rather than the isotropic second-order value
$2^{(d+4)J}/n$.

\section{Elliptic regularity and linearisation}

We collect the PDE estimates used below. All constants are uniform over
$\Fcal(\alpha,r)$ once $r$ is sufficiently small.

\begin{lemma}[Uniform Schauder estimates]\label{lem:Schauder}
There is $r_0\in(0,1)$ such that, uniformly in $f,h\in\Fcal(\alpha,r)$,
\begin{equation}\label{eq:uniform-u}
\|u_f\|_{C^{\alpha+1}(D)}\leq C
\end{equation}
and
\begin{equation}\label{eq:forward-holder}
\|u_f-u_h\|_{C^{\alpha+1}(D)}
\leq
C\|f-h\|_{C^\alpha(D)}.
\end{equation}
In particular,
\[
\nabla u_f\in C^\alpha(D),\qquad
\Delta u_f\in C^{\alpha-1}(D).
\]
\end{lemma}

\begin{proof}
Because $\alpha>1$, the divergence-form equation can be written as
\[
-f\Delta u_f-\nabla f\cdot\nabla u_f=-1.
\]
The coefficients are uniformly in $C^{\alpha-1}$ and $f$ is uniformly
elliptic. The first estimate is therefore the classical Dirichlet Schauder
estimate. For the difference $\delta=u_f-u_h$ and $a=f-h$,
\begin{equation}\label{eq:difference}
\Div(f\nabla\delta)=-\Div(a\nabla u_h).
\end{equation}
The right side belongs to $C^{\alpha-1}$ with norm bounded by
$C\|a\|_{C^\alpha}$, using \eqref{eq:uniform-u}. The Schauder estimate in
divergence form gives \eqref{eq:forward-holder}. See
\cite[Chapters 6 and 8]{GT}.
\end{proof}

For $h\in\Fcal(\alpha,r)$ define
\[
D_h v=-\Div(h\nabla v),
\]
and let $D_h^{-1}$ denote the inverse with homogeneous Dirichlet boundary
conditions. Define the transport-type operator
\[
A_ha=\Div(a\nabla u_h).
\]
The derivative of $F(h)=u_h$ at $h$ is
\begin{equation}\label{eq:derivative}
T_ha=D_h^{-1}A_ha.
\end{equation}
Indeed, $v=T_ha$ solves
\begin{equation}\label{eq:derivative-equation}
\Div(h\nabla v)=-\Div(a\nabla u_h),
\qquad
v|_{\partial D}=0.
\end{equation}

\begin{lemma}[Quadratic expansion]\label{lem:Taylor}
For $h,h+a\in\Fcal(\alpha,r)$,
\[
u_{h+a}=u_h+T_ha+R_h(a),
\]
where
\[
D_hR_h(a)=\Div(a\nabla(u_{h+a}-u_h)).
\]
Moreover, uniformly in $h$,
\begin{equation}\label{eq:remainder-basic}
\|R_h(a)\|_{H^2(D)}
\leq
C\|\Div(a\nabla(u_{h+a}-u_h))\|_{L^2(D)}.
\end{equation}
\end{lemma}

\begin{proof}
Set $\delta=u_{h+a}-u_h$. From \eqref{eq:difference},
\[
\Div((h+a)\nabla\delta)=-\Div(a\nabla u_h).
\]
Therefore
\[
\Div(h\nabla\delta)
=
-\Div(a\nabla u_h)-\Div(a\nabla\delta).
\]
Subtracting the equation for $T_ha$ in \eqref{eq:derivative-equation} gives
\[
\Div(h\nabla R_h(a))=-\Div(a\nabla\delta),
\]
which is equivalent to the asserted formula. The estimate
\eqref{eq:remainder-basic} is the standard $H^2$-estimate for the uniformly
elliptic operator $D_h$.
\end{proof}

At the constant conductivity,
\begin{equation}\label{eq:u-one}
u_{\bone}(x)=\frac{|x|^2-1}{2d},
\qquad
\nabla u_{\bone}(x)=\frac{x}{d},
\qquad
\Delta u_{\bone}=1.
\end{equation}
By Lemma \ref{lem:Schauder}, after reducing $r_0$ if necessary,
\begin{equation}\label{eq:uniform-geometry}
\left\|\nabla u_h-\frac{x}{d}\right\|_{C^1}
+\left\|\Delta u_h-1\right\|_\infty
\leq Cr
\end{equation}
and $\Delta u_h\geq c>0$.

\section{The sharp anisotropic regularisation}

The upper bound is based on a frequency norm adapted to the characteristic
set of $A_h$. Since
\[
A_ha=a\Delta u_h+\nabla a\cdot\nabla u_h,
\]
its principal symbol is $i\nabla u_h(x)\cdot\xi$, while its zeroth-order
symbol is $\Delta u_h(x)$. Hence
\[
\sigma(A_h)(x,\xi)
=
i\nabla u_h(x)\cdot\xi+\Delta u_h(x)+\text{lower order}.
\]
In view of \eqref{eq:uniform-geometry}, in the high-frequency regime this has
size
\[
\abs{\sigma(A_h)(x,\xi)}
\asymp
1+\abs{\frac{x}{|x|}\cdot\xi}.
\]
Because $D_h^{-1}$ is elliptic of order $-2$, the local gain of $T_h$ at
frequency $\xi$ is
\begin{equation}\label{eq:lambda-symbol}
\lambda(x,\xi)
\asymp
\frac{1+\abs{\frac{x}{|x|}\cdot\xi}}{1+|\xi|^2}.
\end{equation}
Thus:
\begin{itemize}
\item for radial covectors, $\lambda\asymp (1+|\xi|)^{-1}$;
\item for tangential covectors, $\lambda\asymp (1+|\xi|)^{-2}$.
\end{itemize}
The tangential covectors form a codimension-one subset of the cosphere
bundle. This codimension-one reduction is responsible for the $d+3$ in the
rate.

\begin{proposition}[Anisotropic inverse stability and noise estimate]
\label{prop:anisotropic}
There exist finite-rank Hilbert seminorms $N_J$, $J\geq 1$, such that the
following holds uniformly for $h,f\in\Fcal(\alpha,r)$, with $r\leq r_0$:
\begin{equation}\label{eq:nonlinear-stability}
\|f-h\|_{L^2(D)}
\leq
C N_J(u_f-u_h)+C2^{-\alpha J},
\end{equation}
and
\begin{equation}\label{eq:noise}
\E\!\left[N_J(\W)^2\right]\leq C2^{(d+3)J}.
\end{equation}
The seminorm $N_J$ extends continuously to the finite-order distributions
appearing in \eqref{eq:model}.
\end{proposition}

The proof is deferred to Section \ref{sec:microlocal}. We first show how it
implies the upper bound.

\section{Upper bound}

Define
\[
Q_n(f)=N_J(Y-u_f)^2.
\]
Let $\widehat f_n\in\Fcal(\alpha,r)$ be an approximate minimiser satisfying
\[
Q_n(\widehat f_n)\leq \inf_{f\in\Fcal(\alpha,r)}Q_n(f)+n^{-1}.
\]
Existence of an exact minimiser follows from compactness of
$\Fcal(\alpha,r)$ in $C^{\alpha'}$ for $\alpha'<\alpha$ and continuity of
$f\mapsto u_f$, but the approximate-minimiser formulation avoids all
measurability issues.

\begin{proposition}\label{prop:upper}
With the above choice,
\[
\sup_{f_0\in\Fcal(\alpha,r)}
\E_{f_0}\|\widehat f_n-f_0\|_{L^2}^2
\leq
C\left(
\frac{2^{(d+3)J}}n+2^{-2\alpha J}
\right).
\]
\end{proposition}

\begin{proof}
By the basic inequality,
\[
N_J(Y-u_{\widehat f_n})
\leq
N_J(Y-u_{f_0})+n^{-1/2}.
\]
Consequently, by the triangle inequality,
\begin{align*}
N_J(u_{\widehat f_n}-u_{f_0})
&\leq
N_J(u_{\widehat f_n}-Y)+N_J(Y-u_{f_0})\\
&\leq
2N_J(Y-u_{f_0})+n^{-1/2}\\
&=
\frac2{\sqrt n}N_J(\W)+n^{-1/2}.
\end{align*}
Squaring and taking expectations, \eqref{eq:noise} gives
\[
\E_{f_0}N_J(u_{\widehat f_n}-u_{f_0})^2
\leq
\frac{C2^{(d+3)J}}n+\frac Cn.
\]
The claim now follows from \eqref{eq:nonlinear-stability}.
\end{proof}

Choose $J=J_n$ so that
\[
2^{J(2\alpha+d+3)}\asymp \frac n{\log n}.
\]
Then
\[
\frac{2^{(d+3)J}}n
\asymp
\frac{2^{-2\alpha J}}{\log n}
\quad\text{and}\quad
2^{-2\alpha J}
\asymp
\left(\frac{\log n}{n}\right)^{\frac{2\alpha}{2\alpha+d+3}}.
\]
Thus
\[
\sup_{f_0\in\Fcal(\alpha,r)}
\E_{f_0}\|\widehat f_n-f_0\|_{L^2}^2
\leq
C n^{-\kappa}(\log n)^\kappa,
\]
proving the upper inequality in Theorem \ref{thm:main}.

\begin{remark}
If the distribution of singular directions is ignored and one simply uses
an isotropic $H^2$-cutoff, the noise dimension is $2^{dJ}$ and the inverse
amplification is $2^{2J}$. This produces
$n^{-2\alpha/(2\alpha+d+4)}$, which is generally suboptimal. The sharp
norm $N_J$ assigns the two-derivative amplification only to the
codimension-one tangential frequency set.
\end{remark}

\section{Lower bound}

\subsection{Radial base conductivity}

Choose a smooth radial function $\rho_0$ supported in
$(1/4,3/4)$ and a small $\eta_b>0$ such that
\[
b(x)=1+\eta_b\rho_0(|x|)
\]
satisfies
\[
\|b-\bone\|_{C^\alpha}\leq \frac r2.
\]
Since $b$ is radial, the solution $u_b$ is radial. More precisely, if
$u_b(x)=U_b(|x|)$, then
\[
r^{-(d-1)}\frac{\dd}{\dd r}
\left(r^{d-1}b(r)U_b'(r)\right)=1,
\]
and therefore
\begin{equation}\label{eq:radial-gradient}
\nabla u_b(x)
=
\frac{x}{d\,b(|x|)}.
\end{equation}
In particular, $\nabla u_b$ is radial.

\subsection{Angular wavelet packing}

Let $\{\psi_{Jm}:m=1,\ldots,M_J\}$ be real-valued spherical needlets on
$S^{d-1}$ at scale $J$ with disjoint angular supports, frequency $2^J$, and
\[
\|\psi_{Jm}\|_{L^2(S^{d-1})}\asymp 1,
\qquad
M_J\asymp 2^{(d-1)J}.
\]
Such systems are standard; see, for example, \cite{Triebel} for
H\"older--Zygmund characterisations and the standard spherical-needlet
construction.

Fix another smooth radial profile $\rho$ supported in $(1/4,3/4)$ with
$\|\rho\|_{L^2(0,1;r^{d-1}\dd r)}=1$, and define
\[
\Psi_{Jm}(r\omega)=\rho(r)\psi_{Jm}(\omega),
\qquad r\in(0,1),\ \omega\in S^{d-1}.
\]
Then
\[
\|\Psi_{Jm}\|_{L^2(D)}\asymp 1,
\qquad
\|\Psi_{Jm}\|_{C^\alpha(D)}
\asymp
2^{J(\alpha+(d-1)/2)}.
\]

Set
\[
c_J=\eta 2^{-J(\alpha+(d-1)/2)},
\]
where $\eta>0$ is chosen small enough that all functions below belong to
$\Fcal(\alpha,r)$. For
$\vartheta=(\vartheta_1,\ldots,\vartheta_{M_J})\in\{0,1\}^{M_J}$, define
\[
f_\vartheta
=
b+c_J\sum_{m=1}^{M_J}\vartheta_m\Psi_{Jm}.
\]
Because the needlets at scale $J$ have disjoint supports,
\[
\|f_\vartheta-\bone\|_{C^\alpha}
\leq
\|b-\bone\|_{C^\alpha}
+
Cc_J2^{J(\alpha+(d-1)/2)}
\leq r
\]
after reducing $\eta$. Hence $f_\vartheta\in\Fcal(\alpha,r)$.

By the Varshamov--Gilbert lemma \cite[Lemma 2.9]{Tsybakov}, there is
$\mathcal M\subset\{0,1\}^{M_J}$ such that
\[
\log|\mathcal M|\geq cM_J
\]
and
\[
\sum_{m=1}^{M_J}|\vartheta_m-\vartheta'_m|
\geq \frac{M_J}{8},
\qquad
\vartheta\neq\vartheta'\in\mathcal M.
\]
Consequently,
\begin{align}
\|f_\vartheta-f_{\vartheta'}\|_{L^2(D)}^2
&=
c_J^2\sum_m(\vartheta_m-\vartheta'_m)^2
\|\Psi_{Jm}\|_{L^2}^2
\notag\\
&\geq
c\,c_J^2M_J
\geq
c\,2^{-2\alpha J}.
\label{eq:packing-separation}
\end{align}

\subsection{Smallness of the forward differences}

\begin{lemma}\label{lem:angular-smoothing}
Uniformly in $\vartheta$,
\[
\|u_{f_\vartheta}-u_b\|_{L^2(D)}
\leq
C2^{-2J}\|f_\vartheta-b\|_{L^2(D)}
\leq
C2^{-2J}2^{-\alpha J}.
\]
\end{lemma}

\begin{proof}
Write $a=f_\vartheta-b$ and $\delta=u_{f_\vartheta}-u_b$. By
\eqref{eq:difference},
\[
-\Div(f_\vartheta\nabla\delta)
=
\Div(a\nabla u_b).
\]
Using \eqref{eq:radial-gradient},
\[
\nabla u_b(r\omega)=\frac{r}{d\,b(r)}\,\omega.
\]
Since $a(r\omega)=\rho(r)\sum_m\vartheta_m\psi_{Jm}(\omega)$, the vector
$a\nabla u_b$ is radial and
\[
\Div(a\nabla u_b)(r\omega)
=
r^{-(d-1)}
\partial_r\!\left(
r^{d-1}\frac{r}{d\,b(r)}\rho(r)
\right)
\sum_m\vartheta_m\psi_{Jm}(\omega).
\]
There is no angular derivative of $a$. Therefore
\[
\|\Div(a\nabla u_b)\|_{L^2(D)}\leq C\|a\|_{L^2(D)},
\]
and this function is at angular frequency $2^J$. Hence
\[
\|\Div(a\nabla u_b)\|_{H^{-2}(D)}
\leq
C2^{-2J}\|a\|_{L^2(D)}.
\]
Uniform elliptic regularity for the operator $D_{f_\vartheta}$ gives
\[
\|\delta\|_{L^2}
\leq
C\|\Div(a\nabla u_b)\|_{H^{-2}}
\leq
C2^{-2J}\|a\|_{L^2}.
\]
Finally,
\[
\|a\|_{L^2}
=
c_J\left(\sum_m\vartheta_m^2\|\Psi_{Jm}\|_2^2\right)^{1/2}
\leq
Cc_JM_J^{1/2}
\leq
C2^{-\alpha J}.
\]
\end{proof}

\subsection{Kullback--Leibler bound and Fano's inequality}

For model \eqref{eq:model}, the Cameron--Martin formula for Gaussian white
noise gives
\begin{equation}\label{eq:KL}
\KL(P_f\|P_g)
=
\frac n2\|u_f-u_g\|_{L^2(D)}^2.
\end{equation}
See \cite[Chapter 6]{GN}. For $\vartheta,\vartheta'\in\mathcal M$, the
triangle inequality and Lemma \ref{lem:angular-smoothing} yield
\begin{align}
\KL(P_{f_\vartheta}\|P_{f_{\vartheta'}})
&\leq
Cn\left(
\|u_{f_\vartheta}-u_b\|_2^2
+
\|u_{f_{\vartheta'}}-u_b\|_2^2
\right)
\notag\\
&\leq
Cn2^{-J(2\alpha+4)}.
\label{eq:KL-bound}
\end{align}
Fano's lemma \cite[Theorem 2.5]{Tsybakov}, combined with
\eqref{eq:packing-separation}, gives
\[
\inf_{\widehat f_n}
\sup_{f_0\in\Fcal(\alpha,r)}
\E_{f_0}\|\widehat f_n-f_0\|_2^2
\geq
c2^{-2\alpha J}
\]
provided
\[
Cn2^{-J(2\alpha+4)}
\leq
c\log|\mathcal M|
\asymp
2^{(d-1)J}.
\]
Thus it suffices that
\[
2^{J(2\alpha+d+3)}\leq c_0n
\]
for a sufficiently small $c_0$. Taking the largest such integer $J$ gives
\[
2^{-2\alpha J}
\geq
c n^{-\frac{2\alpha}{2\alpha+d+3}}.
\]
This proves the lower bound in Theorem \ref{thm:main}.

\section{Microlocal proof of Proposition \ref{prop:anisotropic}}
\label{sec:microlocal}

It remains to prove the anisotropic inverse-stability estimate
\eqref{eq:nonlinear-stability} and the noise estimate \eqref{eq:noise}.

\subsection{Wave-packet decomposition}

For integers $0\leq k\leq j$, consider covectors satisfying
\[
|\xi|\asymp 2^j,
\qquad
\left|\frac{x}{|x|}\cdot\xi\right|\asymp 2^k.
\]
Here $k=0$ means $\left||x|^{-1}x\cdot\xi\right|\lesssim 1$. At fixed $x$,
the angular constraint removes one degree of freedom. A standard
phase-space volume calculation gives the number of coherent states in the
$(j,k)$ cell as
\begin{equation}\label{eq:multiplicity}
M_{jk}\asymp 2^{(d-1)j+k}.
\end{equation}
Indeed, the tangential angular variables contribute $2^{(d-1)j}$ and the
radial frequency contributes $2^k$.

By the standard anisotropic wave-packet construction associated with a
pseudodifferential symbol, there is a frame
\[
\{\psi_{jk\nu}:0\leq k\leq j,\ \nu=1,\ldots,M_{jk}\}
\]
with the following properties:
\begin{enumerate}
\item frame equivalence in $L^2(D)$;
\item frequency localisation to the $(j,k)$ cell;
\item H\"older--Zygmund approximation
\begin{equation}\label{eq:frame-tail}
\left\|
a-\sum_{j\leq J}\sum_{k\leq j}\sum_{\nu}
\langle a,\psi_{jk\nu}\rangle\psi_{jk\nu}
\right\|_{L^2}
\leq
C2^{-\alpha J}\|a\|_{C^\alpha};
\end{equation}
\item on the $(j,k)$ cell, $T_h$ has gain
\begin{equation}\label{eq:local-gain}
\lambda_{jk}\asymp 2^{k-2j}.
\end{equation}
\end{enumerate}
For constant coefficients this is just a local Fourier decomposition. For
$h\in C^\alpha$, $\alpha>1$, the same estimates follow from paraproduct
calculus: the commutator of $A_h$ or the elliptic parametrix for $D_h$ with
a wave-packet cutoff gains one derivative and is therefore negligible on the
H\"older scale. The required paraproduct and symbol-calculus estimates can
be found in \cite[Chapter 3]{TaylorTools}; see also
\cite[Chapter XVIII]{HormanderIII}.

Define
\begin{equation}\label{eq:NJ}
N_J(w)^2
=
\sum_{j\leq J}\sum_{k\leq j}
2^{4j-2k}
\sum_{\nu=1}^{M_{jk}}
|\langle w,\psi_{jk\nu}\rangle|^2.
\end{equation}
Since the sum is finite, $N_J$ is well defined for every finite-order
distribution.

\subsection{Noise estimate}

White noise satisfies
\[
\E|\langle\W,\psi_{jk\nu}\rangle|^2
=
\|\psi_{jk\nu}\|_{L^2}^2
\asymp 1.
\]
Using \eqref{eq:NJ} and \eqref{eq:multiplicity},
\begin{align*}
\E N_J(\W)^2
&\leq
C\sum_{j\leq J}\sum_{k\leq j}
2^{4j-2k}M_{jk}\\
&\leq
C\sum_{j\leq J}\sum_{k\leq j}
2^{4j-2k}2^{(d-1)j+k}\\
&=
C\sum_{j\leq J}\sum_{k\leq j}
2^{(d+3)j-k}\\
&\leq
C2^{(d+3)J}.
\end{align*}
This proves \eqref{eq:noise}.

\subsection{Linear regularised stability}

\begin{lemma}\label{lem:linear-N-stability}
Uniformly in $h\in\Fcal(\alpha,r)$ and $a\in C^\alpha(D)$,
\begin{equation}\label{eq:linear-N-stability}
\|a\|_{L^2}
\leq
CN_J(T_ha)+C2^{-\alpha J}\|a\|_{C^\alpha}.
\end{equation}
\end{lemma}

\begin{proof}
Decompose $a$ into wave packets. On the $(j,k)$ cell,
\eqref{eq:local-gain} says that $T_h$ multiplies the packet coefficient by
$\lambda_{jk}\asymp 2^{k-2j}$, up to a lower-order error. Hence the weight
$2^{2j-k}=\lambda_{jk}^{-1}$ in $N_J$ recovers the packet coefficient from
the corresponding coefficient of $T_ha$. Summing over all cells with
$j\leq J$ gives
\[
\|a_{\leq J}\|_{L^2}
\leq
CN_J(T_ha)+C2^{-J}\|a\|_{C^\alpha},
\]
where the last term comes from commutators and the approximate-parametrix
error. The high-frequency remainder satisfies \eqref{eq:frame-tail}.
Combining these bounds proves \eqref{eq:linear-N-stability}.
\end{proof}

The important point in \eqref{eq:local-gain} is that $k$ measures only the
radial derivative. On the characteristic cells $k=O(1)$, $A_h$ has order
zero and the net gain is $2^{-2j}$; on generic cells $k\asymp j$, the gain
is $2^{-j}$.

\subsection{Quadratic remainder in the $N_J$-norm}

\begin{lemma}\label{lem:remainder-N}
For $r_0$ sufficiently small, uniformly in $h,h+a\in\Fcal(\alpha,r)$,
\begin{equation}\label{eq:remainder-N}
N_J(R_h(a))
\leq
Cr\|a\|_{L^2}
+
C2^{-\alpha J}\|a\|_{C^\alpha}.
\end{equation}
\end{lemma}

\begin{proof}
By Lemma \ref{lem:Taylor},
\[
R_h(a)=D_h^{-1}\Div(a\nabla\delta),
\qquad
\delta=T_ha+R_h(a).
\]
We paralinearise the bilinear term as
\[
\Div(a\nabla T_ha)
=
\Pi_{\nabla T_ha}\cdot\nabla a
+
\Pi_a\Delta T_ha
+
\Pi_{\nabla a}\cdot\nabla T_ha
+
\mathfrak r_h(a),
\]
where $\Pi_b$ denotes Bony's paraproduct. On an anisotropic $(j,k)$ cell,
$A_h$ has order $1$ in the radial variable and order $0$ tangentially, while
$D_h^{-1}$ has order $-2$. Hence the quadratic operator
\[
(a,b)\mapsto D_h^{-1}\Div(a\nabla T_hb)
\]
has one fewer order than the product of the corresponding isotropic
derivatives on every cell. The bilinear paraproduct estimates in
\cite[Chapter 3]{TaylorTools} therefore give
\[
N_J(D_h^{-1}\Div(a\nabla T_ha))
\leq
C\|a\|_{C^\alpha}\|a\|_{L^2}
+
C2^{-\alpha J}\|a\|_{C^\alpha}.
\]
For the term containing $R_h(a)$, the same paraproduct estimate and the
uniform ellipticity of $D_h$ give
\[
N_J(D_h^{-1}\Div(a\nabla R_h(a)))
\leq
CrN_J(R_h(a))
+
C2^{-\alpha J}\|a\|_{C^\alpha}.
\]
Thus
\[
N_J(R_h(a))
\leq
CrN_J(R_h(a))
+
Cr\|a\|_{L^2}
+
C2^{-\alpha J}\|a\|_{C^\alpha}.
\]
After reducing $r_0$ so that $Cr_0<1/2$, the first term on the right is
absorbed into the left side. This proves \eqref{eq:remainder-N}.
\end{proof}

\subsection{Completion of the proof of Proposition
\ref{prop:anisotropic}}

Let $f=h+a$. By Lemmas \ref{lem:linear-N-stability},
\ref{lem:remainder-N}, and the triangle inequality,
\begin{align*}
\|a\|_{L^2}
&\leq
CN_J(T_ha)+C2^{-\alpha J}\|a\|_{C^\alpha}\\
&\leq
CN_J(u_f-u_h)
+CN_J(R_h(a))
+C2^{-\alpha J}\|a\|_{C^\alpha}\\
&\leq
CN_J(u_f-u_h)
+Cr\|a\|_{L^2}
+C2^{-\alpha J}\|a\|_{C^\alpha}.
\end{align*}
Since $\|a\|_{C^\alpha}\leq 2r$, after reducing $r_0$ so that $Cr_0<1/2$,
we obtain
\[
\|a\|_{L^2}
\leq
CN_J(u_f-u_h)+C2^{-\alpha J}.
\]
This is \eqref{eq:nonlinear-stability}. The proof of Proposition
\ref{prop:anisotropic}, and hence of Theorem \ref{thm:main}, is complete.
\qed

\section{Final comparison of exponents}

The sharp exponent is
\[
\boxed{\frac{2\alpha}{2\alpha+d+3}}.
\]
The lower bound comes from angular wave packets at frequency $2^J$:
there are $2^{(d-1)J}$ such modes, and each forward response is of size
$2^{-2J}$. The matching upper bound uses the fact that all other
phase-space directions have a one-derivative gain and therefore contribute
less to the total variance. Treating the inverse problem as uniformly
two-smoothing instead yields only
\[
\frac{2\alpha}{2\alpha+d+4},
\]
which is not the sharp minimax exponent for the full isotropic H\"older
class.

\begin{thebibliography}{99}

\bibitem{GN}
E. Gin\'e and R. Nickl,
\emph{Mathematical Foundations of Infinite-Dimensional Statistical Models}.
Cambridge Series in Statistical and Probabilistic Mathematics,
Cambridge University Press, Cambridge, 2016.

\bibitem{GT}
D. Gilbarg and N. S. Trudinger,
\emph{Elliptic Partial Differential Equations of Second Order}.
Classics in Mathematics, Springer-Verlag, Berlin, 2001.

\bibitem{HormanderIII}
L. H\"ormander,
\emph{The Analysis of Linear Partial Differential Operators III:
Pseudo-Differential Operators}.
Grundlehren der mathematischen Wissenschaften 274,
Springer-Verlag, Berlin, 1985.

\bibitem{TaylorTools}
M. E. Taylor,
\emph{Tools for PDE: Pseudodifferential Operators, Paradifferential
Operators, and Layer Potentials}.
Mathematical Surveys and Monographs 81,
American Mathematical Society, Providence, RI, 2000.

\bibitem{Triebel}
H. Triebel,
\emph{Theory of Function Spaces}.
Monographs in Mathematics 78, Birkh\"auser Verlag, Basel, 1983.

\bibitem{Tsybakov}
A. B. Tsybakov,
\emph{Introduction to Nonparametric Estimation}.
Springer Series in Statistics, Springer, New York, 2009.

\end{thebibliography}

\end{document}
```